% 太阳系（综述）
% license CCBYSA3
% type Wiki

本文根据 CC-BY-SA 协议转载翻译自维基百科\href{https://en.wikipedia.org/wiki/Solar_System}{相关文章}。

\begin{figure}[ht]
\centering
\includegraphics[width=8cm]{./figures/5a727e2db62d7a0e.png}
\caption{太阳、行星、卫星与矮行星[a] （真实色彩，比例为**尺寸等比**，但**距离不按比例**）
```
} \label{fig_Solar_1}
\end{figure}
太阳系（The Solar System）由太阳及其所环绕的天体组成。该名称来源于“Sōl”，即太阳的拉丁名称。大约在 46 亿年前，一个分子云中密度较高的区域发生坍缩，形成了太阳以及一个原行星盘（protoplanetary disc），环绕天体便由此汇聚而成。太阳核心内部的氢融合为氦，释放出能量，这些能量主要通过其外层光球（photosphere）向外辐射，从而在整个系统中形成一个由内而外逐渐降低的温度梯度。太阳系中超过 99.86\% 的质量集中在太阳内部。

在绕太阳运行的天体中，质量最大的为八大行星。按与太阳距离递增的顺序，最靠近太阳的四颗类地行星（terrestrial planets）分别是—— 水星（Mercury）、金星（Venus）、地球（Earth）和火星（Mars）。它们构成了太阳系的内行星系统（inner Solar System）。其中，地球与火星是太阳系中仅有的两颗位于太阳宜居带（habitable zone）内的行星——在此范围内，行星表面可存在液态水。

在距离太阳约 5 个天文单位（AU）之外、越过“霜线”（frost line）的位置，分布着四颗更为庞大的行星：两颗气态巨行星（gas giants）—— 木星（Jupiter）与土星（Saturn），以及两颗冰巨行星（ice giants）—— 天王星（Uranus）与 海王星（Neptune）。它们共同构成了太阳系的外行星系统（outer Solar System）。其中，木星与土星占据了太阳系除恒星外几乎 90\% 的全部质量。

太阳系中还存在着数量庞大但质量较小的天体。天文学家普遍认为，太阳系至少拥有九颗矮行星（dwarf planets），分别是：谷神星（Ceres）、奥库斯（Orcus）、冥王星（Pluto）、哈梅亚（Haumea）、夸奥尔（Quaoar）、鸟神星（Makemake）、共工（Gonggong）、阋神星（Eris）与赛德娜（Sedna）。在这些天体中，有六颗行星、七颗矮行星及若干其他天体拥有自然卫星（natural satellites），即通常所称的“月亮（moons）”。这些卫星的大小范围极广：从如地球之月那样接近矮行星尺度的巨型卫星，到仅有极小质量的“月粒（moonlets）”。此外，太阳系中还存在大量小型天体（small Solar System bodies），包括小行星（asteroids）、彗星（comets）、半人马天体（centaurs）、流星体（meteoroids）以及星际尘埃云（interplanetary dust clouds）。其中一部分天体分布在小行星带（asteroid belt）（位于火星与木星轨道之间）以及柯伊伯带（Kuiper belt）（位于海王星轨道之外）。

在这些天体之间，存在着充满尘埃与粒子的行星际介质（interplanetary medium）。太阳不断释放的带电粒子形成太阳风（solar wind），它向外扩散，构成一个被称为日球层（heliosphere）的巨大空间。在距离太阳约 70–90 个天文单位（AU）处，太阳风受到星际介质（interstellar medium）的阻挡，形成边界——日球顶（heliopause）。这一区域标志着太阳系与星际空间的分界。更远处、约 2,000 个天文单位（AU）以外，是太阳系最外层的假想区域——奥尔特云（Oort cloud）。它被认为是长周期彗星（long-period comets）的来源，其范围可能延伸至太阳系的引力界限（Hill sphere）边缘，约178,000–227,000 AU（2.81–3.59 光年），在此处太阳系的引力势与银河系的引力势达到平衡。目前，太阳系正穿越一个被称为本地星际云（Local Cloud）的星际介质区域。距离太阳系最近的恒星是比邻星（Proxima Centauri），距离约269,000 AU（4.25 光年）。太阳系与比邻星都位于本地泡（Local Bubble）内，这是银河系中一个直径约 1,000 光年（ly）的相对稀薄的区域。

